% RESEARCH_DRAFT_SHA256: 65dd9f6af960f9661355e7471cc5b8c02b9e8ab5d92deb84e0704d08b02479b7
% Layout and source-reference wrappers expanded for the web renderer.
\documentclass[11pt]{article}
\usepackage[margin=25mm]{geometry}
\usepackage{amsthm}
\usepackage{fontspec,unicode-math}
\setmainfont{Latin Modern Roman}
\setmathfont{Latin Modern Math}
\usepackage{amsmath,amssymb,mathtools}
\usepackage{xurl,hyperref}
\hypersetup{colorlinks=true,urlcolor=blue,linkcolor=blue}
\setcounter{secnumdepth}{0}
\setlength{\parindent}{0pt}
\setlength{\parskip}{5pt}
\setlength{\emergencystretch}{3em}
\newcommand{\sref}[2]{\hyperlink{source:#1:#2}{[S#2]}}
\newcommand{\cataloguescope}[1]{\paragraph{Source question.}#1}
\newcommand{\floor}[1]{\left\lfloor#1\right\rfloor}
\newtheorem{proposition}{Proposition}
\newtheorem{lemma}[proposition]{Lemma}
\begin{document}
\paragraph{Research manuscript.} OpenAI \texttt{gpt-6-astra}, reasoning \texttt{ultra}; 11 October 2026.
\paragraph{Preserved source body.} The complete catalogue mathematical text and sources follow with layout and source links expanded, including source qualifications and its historical preparation notes. New research begins at the section headed Research attempt.

\section{Semi-Supervised Sparse Gaussian Classification}\label{statement}

\subsection{Definitions and mathematical statement}

\paragraph{Shared notation.}Standard number systems are denoted by $\mathbb N,\mathbb Z,\mathbb Q,\mathbb R,\mathbb C$; index ranges, computational models and specialized notation follow the cited problem source and the definitions below.


A sparse Gaussian classification model uses Gaussian class-conditional observations with a sparse discriminating parameter. Labeled observations reveal class membership; unlabeled observations reveal only the feature vectors. Statistical benefit is a decrease in the labeled sample requirement or achievable risk.

A learning problem specifies a data distribution, a hypothesis class and a loss. Statistical consistency means that the estimator approaches its population target as the sample size increases; statistical and computational requirements are distinct.
\paragraph{Statement recorded in the supplied source.}
Characterize the statistical and computational regimes in which unlabeled data provably improve sparse Gaussian classification. \href{https://doi.org/10.4171/owr/2024/30}{[S2]}
\paragraph{Scope qualification.}
The covariance, class proportions, sparsity scaling, and loss are not specified in the short question.
\paragraph{Source status.}
The supplied catalogue labels this record open. Its source-status metadata was checked against the supplied snapshot on 2026-10-11. This import does not independently establish current open status or mathematical correctness. The catalogue literature assessment is dated 2026-08-22.
\subsection{Short English statement}
Characterize when unlabeled examples improve sparse Gaussian classification, accounting for both statistics and computation.
\subsection{Sources}
\begin{itemize}
\item[S1] Ulam AI and the UnsolvedMath Contributors, UnsolvedMath: A Curated Collection of Open Mathematics Problems, catalogue entry. Catalogue curation and metadata: CC BY 4.0. The supplied source question is retained; definitions and the English summary are editorial additions.
\url{https://huggingface.co/datasets/ulamai/UnsolvedMath}.
\item[S2] Statistics and Learning Theory in the Era of Artificial Intelligence (2024), Oberwolfach Reports, DOI 10.4171/owr/2024/30, pp. 1677--1718
\url{https://doi.org/10.4171/owr/2024/30}.
\item[S3] Semi-Supervised Sparse Gaussian Classification
\url{https://arxiv.org/abs/2409.03335}.
\textit{Underlying Oberwolfach Reports material retains its source-specific licence, including CC BY-SA 4.0 where applicable; see THIRD\_PARTY\_NOTICES.md.}
\end{itemize}
\paragraph{Preparation and review.}
Original question attributed to the cited source and UnsolvedMath contributors. Codex (GPT-6; exact serving revision not exposed) assisted with definitions, English summaries, source inspection and formatting. No new human mathematical authorship or independent proof verification is asserted. Catalogue identity checked against the complete supplied registry and existing repository statements; known flagged repetitions and obvious identical conjectures excluded. Full mathematical-equivalence/literature review has not been performed. Statement preparation and final review followed contribution prompts C and D; automated checks verify structure and identity, not open status.
\label{end:statement}

\section{Research attempt}
\subsection{Literature, source context and scope}
The full original catalogue body and source qualifications are preserved above. The relevant primary workshop passages were read in the cached publisher report. Historical claims and subsequent attributed results are distinguished from the mathematical work below. A restricted example, standard lemma or known theorem is not counted as new progress on the entire original target. Missing assumptions and repeated or rhetorical extracts are not silently turned into different solved problems.

The workshop's sparse Gaussian passage reports results, rather than specifying a complete new theorem. Azar and Nadler's primary abstract reports a favorable joint-data regime and conditional computational lower bounds; those claims are context, not results established here. The missing covariance, loss, scaling and hardness assumptions prevent identifying the reconstructed broad question with a single formal target.
\subsection{Attempt: separate support recovery from label orientation}
Fix integers $d,N,m\ge1$, a known $a>0$, and the explicitly restricted model $Y\in\{-1,1\}$ uniform, $X=Y\mu+Z$, $Z\sim N(0,I_d)$ independent, $\mu=sa e_J$ with $s\in\{-1,1\}$. Unlabeled and labeled samples are independent. This is a one-sparse illustration, not the full paper's regime.
\begin{proposition}
Unlabeled observations alone cannot uniformly achieve classification error below $1/2$ over the two signs at a fixed $J$. A polynomial-time two-stage procedure using both sample types has expected excess risk over the Bayes error $r_*=\Phi(-a)$ at most
\[
 \min\{1,(8d+16a^2)/(Na^4)\}+(1-2r_*)\Phi(-a\sqrt m).
\]
\end{proposition}
\begin{proof}
The marginal law of $X$ is the same for $\mu$ and $-\mu$, while their joint laws exchange the labels. For any classifier, including one randomized from an unlabeled training sample, the two expected risks therefore sum to one: couple the training observations, test $X$, and classifier coins identically, and reverse the test label. This proves the first assertion.

For each coordinate compute $V_j=N^{-1}\sum_{i=1}^N X_{ij}^2$ and choose an index maximizing $V_j$. Off the signal coordinate the mean and variance of $X_j^2$ are $1,2$; on the signal they are $1+a^2,2+4a^2$. These follow by expanding $(Z\pm a)^2$ and using $E Z^2=1$, $E Z^4=3$, and vanishing odd moments. Chebyshev's inequality and a union bound show that all empirical moments differ from their means by less than $a^2/2$, except with probability at most $(8d+16a^2)/(Na^4)$. On that event's complement the selected index is $J$ (boundary ties have probability zero).

Using the independent labeled data, estimate the sign by that of $\sum_{i=1}^mY_iX_{iJ}$. This sum has law $N(msa,m)$, hence the wrong-sign probability is $\Phi(-a\sqrt m)$. The correct and reversed threshold classifiers have risks $r_*$ and $1-r_*$, respectively; this follows directly from $YX_J\sim N(sa,1)$. A wrong support choice costs at most one in excess risk. Conditional independence of the two samples now gives the displayed bound. Computing all moments and the labeled sum takes $O(Nd+m)$ arithmetic operations after selection.
\end{proof}
\subsection{Failed inference and original scope}
The orientation argument does not give a labeled-only computational lower bound, and support-recovery error is not automatically a classification lower bound. In particular, the crude Chebyshev support cost is not a sharp statistical-computational transition. Thus the attempted route isolates a real benefit of combining information types but does not reproduce or improve the broad claimed optimal regime or its conditional hardness separation.
\section{Verification and remaining gaps}
Finite checks evaluate the analytic risk bound and the Gaussian moments. The proof uses explicit model assumptions absent from the original extract; unrestricted sparsity, covariance, optimal sample tradeoffs and computational lower bounds remain unproved here.

\section{References and source attribution}
Primary workshop report: \url{https://doi.org/10.4171/owr/2024/30}. Relevant passages read: OWR 2024/30, printed pages 1690--1691, sparse Gaussian classification.

\url{https://arxiv.org/abs/2409.03335}. Read author abstract: reports SSL regime and feature-selection context; full paper proof not audited.

All retrievals were checked on 11 October 2026. Original question attribution remains in the preserved Sources section. External theorems are attributed and not claimed as new discoveries or fully verified merely by citation.


\paragraph{Integration review.} Codex (GPT-6; exact serving revision not exposed) reviewed the written arguments and expanded layout and source-reference wrappers. Repository metadata and the production web renderer were checked using Python 3.13 and JavaScript with MathJax 3.2.2. This model-assisted review is not independent human proof verification; no PDF compilation or formal proof checking is asserted.

\section{Completion Estimate}
\noindent\textbf{COMPLETION: 0\%}
No substantiated new progress toward the entire original target is claimed. The worked results and failed route are delimited above; a known result, illustrative model, or auxiliary obstruction is not counted as original-target completion.

\section{Final status and attribution}
\textbf{Attempt status: unresolved.} This describes this research attempt, not a new ruling on the source catalogue's status or on external literature claims.
The mathematics and this self-audit were generated by OpenAI \texttt{gpt-6-astra}, reasoning \texttt{ultra}, on 11 October 2026. Codex (GPT-6; exact serving revision not exposed) orchestrated the batch, read and reviewed the written mathematical arguments, independently reran the finite checks, and contributed formatting/repository review. No human mathematical authorship, independent human proof checking, or formal verification is claimed. The historical preparation paragraphs in the reproduced catalogue body concern that body only.

Contribution prompt A was applied to the complete definition and supplied source record, followed by the adversarial checks of prompt B. Local reproduction files are \texttt{batch026/checks.py}, \texttt{computations.json}, and the per-ID audit JSON. They were not public links at preparation time. Numerical checks are finite and do not prove unrestricted results. No LaTeX engine was found locally; compilation was not performed by this research agent.

\end{document}
